% --- Archon physics formalization source begin ---
% archon:physics
% archon:covers PhyXMiniProblems/problem_phyx_mini_0079.lean
% archon:source-report reports/phyx_mini/problem_phyx_mini_0079.source.json

\chapter{Physics problem phyx\_mini\_0079}
\label{ch:PhyXMiniProblems_problem_phyx_mini_0079}

\paragraph{Problem source.}
\#\# Physical scenario

In figure, let a beam of x rays of wavelength 0.125\textbackslash{},\textbackslash{}text\{nm\} be incident on an NaCl crystal at angle \textbackslash{}( 	heta = 45.0\textasciicircum{}\textbackslash{}circ \textbackslash{}) to the top face of the crystal and a family of reflecting planes. Let the reflecting planes have separation \textbackslash{}( d = 0.252\textbackslash{},\textbackslash{}text\{nm\} \textbackslash{}). The crystal is turned through angle \textbackslash{}( \textbackslash{}phi \textbackslash{}) around an axis perpendicular to the plane of the page until these reflecting planes give diffraction maxima.

\#\# Question

What is the larger value of \textbackslash{}( \textbackslash{}phi \textbackslash{}) if it is turned counterclockwise?

\#\# Answer choices

- A: A: 31.0\textasciicircum{}\textbackslash{}circ

- B: B: 48.0\textasciicircum{}\textbackslash{}circ

- C: C: 37.8\textasciicircum{}\textbackslash{}circ

- D: D: 41.4\textasciicircum{}\textbackslash{}circ

\#\# Figure caption (auxiliary; use the image as primary evidence)

The image depicts a diagram related to the reflection or diffraction of a beam. 

- There is a horizontal thick black line at the top, representing a surface where an "Incident beam" strikes.
- A red arrow labeled "Incident beam" approaches the surface at an angle \textbackslash{}( \textbackslash{}theta \textbackslash{}) to the normal (represented by a dashed vertical line).
- The angle \textbackslash{}( \textbackslash{}theta \textbackslash{}) is marked between the incident beam and the normal to the surface.
- Below the top surface, there are two more parallel horizontal lines, possibly indicating additional surfaces or layers.
- The distance between these parallel lines is labeled \textbackslash{}( d \textbackslash{}) on the right side, indicating equal spacing between them. 

This setup suggests a scenario involving reflection or diffraction, common in optics and physics contexts.

\#\# Dataset metadata

Wave Optics; Physical Model Grounding Reasoning, Multi-Formula Reasoning

\paragraph{Recorded answer/context.}
C: C: 37.8\textasciicircum{}\textbackslash{}circ

\paragraph{Figure/image path.}
/root/proposal\_for\_physic/science-mango/phyx\_mini\_run/phyx\_data/test\_image/79.png

\paragraph{Formalization target.}
create a compiling Lean file with sorry bodies at `PhyXMiniProblems/problem\_phyx\_mini\_0079.lean`. The Lean declarations must preserve the physical quantities, dimensions or dimensional roles, figure labels, governing-law hypotheses, and final relation expressed by this problem.
Use Mathlib/Physlib names found through LeanExplore where available. If a physics API is missing, introduce faithful local abstractions rather than scalar placeholder aliases.

\begin{theorem}[Physics formalization target]
\label{thm:physics:phyx_mini_0079:target}
\lean{PhyXMiniProblems.ProblemPhyXMini0079.problem_phyx_mini_0079}
\uses{def:physics:phyx-mini-0079:phyxminiproblems-problemphyxmini0079-xraycrystaldiffractionsetup, def:physics:phyx-mini-0079:phyxminiproblems-problemphyxmini0079-matchesproblemandfigurereadouts, def:physics:phyx-mini-0079:phyxminiproblems-problemphyxmini0079-satisfiescounterclockwiserotationgeometry, def:physics:phyx-mini-0079:phyxminiproblems-problemphyxmini0079-satisfiesidealbraggdiffractionlaw, def:physics:phyx-mini-0079:phyxminiproblems-problemphyxmini0079-islargercounterclockwisemaximum, def:physics:phyx-mini-0079:phyxminiproblems-problemphyxmini0079-roundstonearesttenthdegree, def:physics:phyx-mini-0079:phyxminiproblems-problemphyxmini0079-matchesanswer}
The assigned autoformalize agent should translate this physics problem into Lean declarations in the covered file, with theorem and lemma proof bodies written as `by sorry`.
\end{theorem}
\begin{proof}
This is an autoformalization task, not a proof task. Produce statements that can later be proved by the physics prover without weakening the physical model.
\end{proof}

% --- Archon declaration topology begin ---
\section{Declaration topology}
\begin{definition}[Length Quantity]
  \label{def:physics:phyx-mini-0079:phyxminiproblems-problemphyxmini0079-lengthquantity}
  \lean{PhyXMiniProblems.ProblemPhyXMini0079.LengthQuantity}
  A physical length independent of the unit chosen to read it
\end{definition}
\begin{proof}
  This is the declared model definition of Length Quantity.
\end{proof}
\begin{definition}[nanometer Unit Choices]
  \label{def:physics:phyx-mini-0079:phyxminiproblems-problemphyxmini0079-nanometerunitchoices}
  \lean{PhyXMiniProblems.ProblemPhyXMini0079.nanometerUnitChoices}
  Unit choices whose length unit is the nanometer used in the problem
\end{definition}
\begin{proof}
  This is the declared model definition of nanometer Unit Choices.
\end{proof}
\begin{definition}[length In Nanometers]
  \label{def:physics:phyx-mini-0079:phyxminiproblems-problemphyxmini0079-lengthinnanometers}
  \lean{PhyXMiniProblems.ProblemPhyXMini0079.lengthInNanometers}
  \uses{def:physics:phyx-mini-0079:phyxminiproblems-problemphyxmini0079-lengthquantity, def:physics:phyx-mini-0079:phyxminiproblems-problemphyxmini0079-nanometerunitchoices}
  The scalar nanometer readout of a physical length
\end{definition}
\begin{proof}
  This is the declared model definition of length In Nanometers.
\end{proof}
\begin{definition}[degrees To Radians]
  \label{def:physics:phyx-mini-0079:phyxminiproblems-problemphyxmini0079-degreestoradians}
  \lean{PhyXMiniProblems.ProblemPhyXMini0079.degreesToRadians}
  Convert a degree readout into the radian scalar expected by Real.sin
\end{definition}
\begin{proof}
  This is the declared model definition of degrees To Radians.
\end{proof}
\begin{definition}[radians To Degrees]
  \label{def:physics:phyx-mini-0079:phyxminiproblems-problemphyxmini0079-radianstodegrees}
  \lean{PhyXMiniProblems.ProblemPhyXMini0079.radiansToDegrees}
  Convert a radian scalar into a degree readout
\end{definition}
\begin{proof}
  This is the declared model definition of radians To Degrees.
\end{proof}
\begin{definition}[Radiation Kind]
  \label{def:physics:phyx-mini-0079:phyxminiproblems-problemphyxmini0079-radiationkind}
  \lean{PhyXMiniProblems.ProblemPhyXMini0079.RadiationKind}
  The kind of incident radiation distinguished in this problem
\end{definition}
\begin{proof}
  This is the declared model definition of Radiation Kind.
\end{proof}
\begin{definition}[Crystal Material]
  \label{def:physics:phyx-mini-0079:phyxminiproblems-problemphyxmini0079-crystalmaterial}
  \lean{PhyXMiniProblems.ProblemPhyXMini0079.CrystalMaterial}
  The crystal material named in the problem
\end{definition}
\begin{proof}
  This is the declared model definition of Crystal Material.
\end{proof}
\begin{definition}[Rotation Axis Orientation]
  \label{def:physics:phyx-mini-0079:phyxminiproblems-problemphyxmini0079-rotationaxisorientation}
  \lean{PhyXMiniProblems.ProblemPhyXMini0079.RotationAxisOrientation}
  The orientation of the axis about which the crystal is turned
\end{definition}
\begin{proof}
  This is the declared model definition of Rotation Axis Orientation.
\end{proof}
\begin{definition}[Crystal Plane Label]
  \label{def:physics:phyx-mini-0079:phyxminiproblems-problemphyxmini0079-crystalplanelabel}
  \lean{PhyXMiniProblems.ProblemPhyXMini0079.CrystalPlaneLabel}
  The three parallel horizontal lines explicitly drawn in the primary figure
\end{definition}
\begin{proof}
  This is the declared model definition of Crystal Plane Label.
\end{proof}
\begin{definition}[XRay Crystal Diffraction Setup]
  \label{def:physics:phyx-mini-0079:phyxminiproblems-problemphyxmini0079-xraycrystaldiffractionsetup}
  \lean{PhyXMiniProblems.ProblemPhyXMini0079.XRayCrystalDiffractionSetup}
  \uses{def:physics:phyx-mini-0079:phyxminiproblems-problemphyxmini0079-lengthquantity, def:physics:phyx-mini-0079:phyxminiproblems-problemphyxmini0079-radiationkind, def:physics:phyx-mini-0079:phyxminiproblems-problemphyxmini0079-crystalmaterial, def:physics:phyx-mini-0079:phyxminiproblems-problemphyxmini0079-rotationaxisorientation, def:physics:phyx-mini-0079:phyxminiproblems-problemphyxmini0079-crystalplanelabel}
  Physical quantities, labeled figure geometry, and the experimental maximum predicate for the rotating-crystal setup. rotationRadians is positive for counterclockwise rotation. The function glancingAngleAfterCounterclockwiseRotation records the acute angle between the incident beam and the rotated reflecting planes; it is not defined here in terms of the unknown answer
\end{definition}
\begin{proof}
  This is the declared model definition of XRay Crystal Diffraction Setup.
\end{proof}
\begin{definition}[Matches Problem And Figure Readouts]
  \label{def:physics:phyx-mini-0079:phyxminiproblems-problemphyxmini0079-matchesproblemandfigurereadouts}
  \lean{PhyXMiniProblems.ProblemPhyXMini0079.MatchesProblemAndFigureReadouts}
  \uses{def:physics:phyx-mini-0079:phyxminiproblems-problemphyxmini0079-lengthinnanometers, def:physics:phyx-mini-0079:phyxminiproblems-problemphyxmini0079-degreestoradians, def:physics:phyx-mini-0079:phyxminiproblems-problemphyxmini0079-xraycrystaldiffractionsetup}
  The problem-statement and primary-figure readouts. The primary image shows θ between the incident ray and the horizontal top face, and all three drawn planes are parallel. No value of the requested rotation occurs here
\end{definition}
\begin{proof}
  This is the declared model definition of Matches Problem And Figure Readouts.
\end{proof}
\begin{definition}[Satisfies Counterclockwise Rotation Geometry]
  \label{def:physics:phyx-mini-0079:phyxminiproblems-problemphyxmini0079-satisfiescounterclockwiserotationgeometry}
  \lean{PhyXMiniProblems.ProblemPhyXMini0079.SatisfiesCounterclockwiseRotationGeometry}
  \uses{def:physics:phyx-mini-0079:phyxminiproblems-problemphyxmini0079-xraycrystaldiffractionsetup}
  Counterclockwise rotation through φ increases the glancing angle from θ to θ + φ, as read from the primary figure
\end{definition}
\begin{proof}
  This is the declared model definition of Satisfies Counterclockwise Rotation Geometry.
\end{proof}
\begin{definition}[Is Bragg Maximum Of Order]
  \label{def:physics:phyx-mini-0079:phyxminiproblems-problemphyxmini0079-isbraggmaximumoforder}
  \lean{PhyXMiniProblems.ProblemPhyXMini0079.IsBraggMaximumOfOrder}
  \uses{def:physics:phyx-mini-0079:phyxminiproblems-problemphyxmini0079-xraycrystaldiffractionsetup}
  The order-n Bragg relation 2 d sin(α) = n λ, required in every unit system. Both sides are scalar readouts of physical lengths in the same chosen unit
\end{definition}
\begin{proof}
  This is the declared model definition of Is Bragg Maximum Of Order.
\end{proof}
\begin{definition}[Satisfies Ideal Bragg Diffraction Law]
  \label{def:physics:phyx-mini-0079:phyxminiproblems-problemphyxmini0079-satisfiesidealbraggdiffractionlaw}
  \lean{PhyXMiniProblems.ProblemPhyXMini0079.SatisfiesIdealBraggDiffractionLaw}
  \uses{def:physics:phyx-mini-0079:phyxminiproblems-problemphyxmini0079-xraycrystaldiffractionsetup, def:physics:phyx-mini-0079:phyxminiproblems-problemphyxmini0079-isbraggmaximumoforder}
  Ideal Bragg diffraction is the governing physical law: a rotation gives a maximum exactly when the relation holds for some positive integer order
\end{definition}
\begin{proof}
  This is the declared model definition of Satisfies Ideal Bragg Diffraction Law.
\end{proof}
\begin{definition}[Is Admissible Counterclockwise Maximum]
  \label{def:physics:phyx-mini-0079:phyxminiproblems-problemphyxmini0079-isadmissiblecounterclockwisemaximum}
  \lean{PhyXMiniProblems.ProblemPhyXMini0079.IsAdmissibleCounterclockwiseMaximum}
  \uses{def:physics:phyx-mini-0079:phyxminiproblems-problemphyxmini0079-xraycrystaldiffractionsetup}
  A physically admissible counterclockwise solution on the acute glancing-angle branch. This branch contains the positive maxima reached from the pictured 45° starting configuration before the planes become perpendicular to the incident beam
\end{definition}
\begin{proof}
  This is the declared model definition of Is Admissible Counterclockwise Maximum.
\end{proof}
\begin{definition}[Is Larger Counterclockwise Maximum]
  \label{def:physics:phyx-mini-0079:phyxminiproblems-problemphyxmini0079-islargercounterclockwisemaximum}
  \lean{PhyXMiniProblems.ProblemPhyXMini0079.IsLargerCounterclockwiseMaximum}
  \uses{def:physics:phyx-mini-0079:phyxminiproblems-problemphyxmini0079-xraycrystaldiffractionsetup, def:physics:phyx-mini-0079:phyxminiproblems-problemphyxmini0079-isadmissiblecounterclockwisemaximum}
  The selected rotation is the largest admissible counterclockwise maximum
\end{definition}
\begin{proof}
  This is the declared model definition of Is Larger Counterclockwise Maximum.
\end{proof}
\begin{definition}[Answer Choice]
  \label{def:physics:phyx-mini-0079:phyxminiproblems-problemphyxmini0079-answerchoice}
  \lean{PhyXMiniProblems.ProblemPhyXMini0079.AnswerChoice}
  Labels of the four multiple-choice answers
\end{definition}
\begin{proof}
  This is the declared model definition of Answer Choice.
\end{proof}
\begin{definition}[rotation Degrees]
  \label{def:physics:phyx-mini-0079:phyxminiproblems-problemphyxmini0079-answerchoice-rotationdegrees}
  \lean{PhyXMiniProblems.ProblemPhyXMini0079.AnswerChoice.rotationDegrees}
  \uses{def:physics:phyx-mini-0079:phyxminiproblems-problemphyxmini0079-answerchoice}
  The counterclockwise rotation in degrees printed beside each answer
\end{definition}
\begin{proof}
  This is the declared model definition of rotation Degrees.
\end{proof}
\begin{definition}[Rounds To Nearest Tenth Degree]
  \label{def:physics:phyx-mini-0079:phyxminiproblems-problemphyxmini0079-roundstonearesttenthdegree}
  \lean{PhyXMiniProblems.ProblemPhyXMini0079.RoundsToNearestTenthDegree}
  \uses{def:physics:phyx-mini-0079:phyxminiproblems-problemphyxmini0079-radianstodegrees}
  Agreement with a one-decimal-place degree readout, to the nearest tenth
\end{definition}
\begin{proof}
  This is the declared model definition of Rounds To Nearest Tenth Degree.
\end{proof}
\begin{definition}[Matches Answer]
  \label{def:physics:phyx-mini-0079:phyxminiproblems-problemphyxmini0079-matchesanswer}
  \lean{PhyXMiniProblems.ProblemPhyXMini0079.MatchesAnswer}
  \uses{def:physics:phyx-mini-0079:phyxminiproblems-problemphyxmini0079-answerchoice, def:physics:phyx-mini-0079:phyxminiproblems-problemphyxmini0079-roundstonearesttenthdegree}
  A physical rotation rounds to the degree readout of a displayed answer
\end{definition}
\begin{proof}
  This is the declared model definition of Matches Answer.
\end{proof}
\begin{lemma}[larger counterclockwise rotation formula]
  \label{lem:physics:phyx-mini-0079:phyxminiproblems-problemphyxmini0079-larger-counterclockwise-rotation-formula}
  \lean{PhyXMiniProblems.ProblemPhyXMini0079.larger_counterclockwise_rotation_formula}
  \uses{def:physics:phyx-mini-0079:phyxminiproblems-problemphyxmini0079-lengthinnanometers, def:physics:phyx-mini-0079:phyxminiproblems-problemphyxmini0079-xraycrystaldiffractionsetup, def:physics:phyx-mini-0079:phyxminiproblems-problemphyxmini0079-matchesproblemandfigurereadouts, def:physics:phyx-mini-0079:phyxminiproblems-problemphyxmini0079-satisfiescounterclockwiserotationgeometry, def:physics:phyx-mini-0079:phyxminiproblems-problemphyxmini0079-satisfiesidealbraggdiffractionlaw, def:physics:phyx-mini-0079:phyxminiproblems-problemphyxmini0079-islargercounterclockwisemaximum}
  The larger admissible branch is the fourth-order Bragg maximum. Its exact rotation is the fourth-order Bragg angle minus the initial glancing angle. This is a derived intermediate result, not a supplied readout
\end{lemma}
\begin{proof}
  The conclusion follows from the governing relations and figure readouts named in the statement of larger counterclockwise rotation formula.
\end{proof}
% --- Archon declaration topology end ---
% --- Archon physics formalization source end ---
